\chapter{Introduction}
\label{ch:intro}

\section{What \fxme\ is}

\fxme\ lets your digital audio workstation (DAW) run its effects on another
computer. In the DAW you insert one plugin, called \ui{FxmeNet}. Behind it,
a chain of plugins runs on a \emph{server} (the program
\cmd{fxmenet-server}) on the same local network. The audio goes to the
server and comes back; the DAW sees one ordinary plugin, and compensates the
delay it reports, as it does for any plugin with latency.

The reason is simple: a session with many heavy plugins can exhaust the CPU
of the computer that runs the DAW. A second machine, even an old one, can
take that load. In our tests (chapter~\ref{ch:performance}), 24 tracks with
17 plugins each used 75\,\% of the client's processor when processed
locally, and 20 to 25\,\% through \fxme.

The server is a console program. It needs no screen, no desktop session and
no keyboard, so it can live in a cupboard or a rack, and it can run on a
different operating system than the DAW.

\section{How it works}

\fxme\ has two parts.

\begin{itemize}
  \item The \textbf{client} is the \ui{FxmeNet} plugin (VST3 and AU), loaded
    in the DAW. It has the editor you work with.
  \item The \textbf{server} is \cmd{fxmenet-server}. It loads the VST3
    plugins you ask for and processes the audio.
\end{itemize}

\begin{lstlisting}[language={}]
   Your DAW (client)                          Server
 +---------------------+                +-------------------+
 |  track 1: FxmeNet --+--- audio UDP --->  session 1: chain |
 |                     <--- audio UDP ---+  (plugin, plugin) |
 |  track 2: FxmeNet --+--- audio UDP --->  session 2: chain |
 |                     <--- audio UDP ---+                   |
 |  every instance  ---+--- control TCP -->  plugin list,    |
 |                     <-------------------  parameters, ... |
 +---------------------+                +-------------------+
\end{lstlisting}

Each instance of \ui{FxmeNet} in your project opens one \emph{session} on
the server, with its own chain of plugins. Two kinds of traffic flow:

\begin{itemize}
  \item \textbf{Audio}, in small blocks of N samples, over UDP, in both
    directions. A reply that is late is replaced by a short fade to silence
    and back (the block is \emph{concealed}); this is how a network that is
    sometimes slow stays inaudible most of the time.
  \item \textbf{Control}, over TCP: the list of plugins, adding and
    removing plugins, parameters, saving and restoring state.
\end{itemize}

The \emph{network budget} D says how many blocks of time a reply has to come
back. The latency that \fxme\ adds is about N times D samples, plus what the
DAW's own block size adds. Chapter~\ref{ch:performance} explains how to
choose them.

\section{What you see: native editors}

When you work on a plugin that sits on the server, you do not look at a
picture of a remote window. If the same plugin is installed on the client
computer, \fxme\ opens the plugin's own editor locally, and keeps it in
step with the copy on the server. Knobs respond at once, and the editor
scales with your screen. If the plugin is made for \fxme\ (see
chapter~\ref{ch:developers}), its meters and analysers show the server's
real signal. Plugins that are not installed on the client still work, with a
generic panel of sliders. These cases are the \emph{tiers} of
chapter~\ref{ch:tiers}.

\section{What \fxme\ is not}

\begin{itemize}
  \item It is not a way to play live through a very low latency. The added
    latency is a few milliseconds at best (a few tens of milliseconds for
    heavy loads). It suits mixing and production; for live monitoring,
    keep latency-critical plugins local.
  \item It is for a local network you control. The traffic is not
    encrypted and the server does not ask for a password.
    Do not expose its ports to the internet.
  \item The server hosts VST3 \emph{effects}. Instruments are not offered.
  \item It does not make plugins sound different: the server runs the same
    plugin code your DAW would run.
\end{itemize}

\section{What you need}

\begin{description}
  \item[A server.] A computer with spare CPU and a wired network
    connection. Linux is the reference platform (our tests use Ubuntu
    24.04). \todo{state which other server operating systems are supported}
  \item[A client.] A DAW on Windows, macOS or Linux that hosts VST3 (or AU on
    macOS). The tests used Reaper on Windows. \todo{list the DAWs that
    were tried}
  \item[A network.] A wired gigabit network between the two machines.
    Wi-Fi works badly for audio. A stereo track at 44.1\,kHz uses about
    3\,Mbit/s in each direction (chapter~\ref{ch:network}).
  \item[The plugins.] The plugins you want to run must be installed on the
    server. For their own editors on the client, install the same version
    there too.
\end{description}

\section{Terms used in this manual}

\begin{description}
  \item[Block size N] Samples sent to the server at a time (32 to 1024).
  \item[Network budget D] Blocks of time allowed for a reply to come back (1
    to 8).
  \item[Session] One instance of the \ui{FxmeNet} plugin, connected to the
    server, with its chain.
  \item[Chain] The plugins of a session, processed in an order you choose,
    and optionally wired by the routing (chapter~\ref{ch:routing}).
  \item[Concealed block] A block whose reply came too late, replaced by a
    short fade out and in.
\end{description}

A fuller list is in the glossary (chapter~\ref{ch:glossary}).
